\chapter{Riemannian Geometry Background}

\section{Notation and conventions}

Fix vector fields $X,Y,W,Z$ on a manifold $\M$ and tensor fields $A,B$. The
manifold is equipped with a Riemannian metric $g$, or with a smooth family
$g_t$; the metric and its Levi-Civita connection $\nabla$ act on all tensor
powers. Metric compatibility gives the tensor product rule
\[
\nabla_X(A\otimes B)=(\nabla_XA)\otimes B+A\otimes(\nabla_XB).
\]
When the first argument of a covariant derivative is displayed as a slot, we
write $\nabla_XA=\nabla A(X,\ldots)$.

\begin{definition}[Second covariant derivative]
\dcref{ch2:1.1}
\label{topping:notn-second-covariant-derivative}
\source{topping:page-0026}
\group{topping-background}
\lean{Topping.secondCovDerivAlong, Topping.secondCovDeriv,
  Topping.secondCovDeriv_cons}
\leanok
For a tensor field, define
\[
\nabla^2_{X,Y}:=\nabla_X\nabla_Y-\nabla_{\nabla_XY}.
\]
\end{definition}

\begin{definition}[Hessian]
\dcref{ch2:1.1}
\label{topping:notn-hessian}
\source{topping:page-0026}
\group{topping-background}
\lean{Topping.hessian, Topping.hessian_apply}
\leanok
For $f:\M\to\mathbb{R}$, set
\[
\Hess(f):=\nabla df.
\]
\end{definition}

\begin{definition}[Connection Laplacian]
\dcref{ch2:1.1}
\label{topping:notn-rough-laplacian}
\source{topping:page-0026}
\group{topping-background}
\lean{Topping.roughLaplacian, Topping.roughLaplacian_apply,
  Topping.roughLaplacian_scalarTensorField}
\leanok
The connection Laplacian of a tensor field $A$ is
\[
\Delta A:=\tr_{12}\nabla^2A.
\]
\end{definition}

\begin{definition}[Curvature operator]
\dcref{ch2:1.1}
\label{topping:notn-curvature-operator}
\source{topping:page-0026}
\group{topping-background}
\lean{Topping.curvatureOperator, Topping.curvatureOperator_apply}
\leanok
The curvature operator is
\[
R(X,Y):=\nabla_Y\nabla_X-\nabla_X\nabla_Y+\nabla_{[X,Y]}.
\]
\end{definition}

\begin{definition}[Riemann curvature tensor]
\dcref{ch2:1.1}
\label{topping:notn-riemann-tensor}
\source{topping:page-0026}
\group{topping-background}
\lean{Topping.riemannCurvature, Topping.riemannCurvatureAt, Topping.riemannCurvatureAt_eq}
\leanok
Define
\[
\Rm(X,Y,W,Z):=\langle R(X,Y)W,Z\rangle.
\]
\end{definition}

\begin{definition}[Ricci tensor and scalar curvature]
\dcref{ch2:1.1}
\label{topping:notn-richi-tensor}
\source{topping:page-0027}
\group{topping-background}
\lean{Topping.ricciTensorAt, Topping.ricciTensorAt_eq_sum,
  Topping.scalarCurvatureAt, Topping.scalarCurvatureAt_eq_trace}
\leanok
The contractions are
\[
\Ric(X,Y):=\tr\Rm(X,\cdot,Y,\cdot),\qquad R:=\tr\Ric.
\]
\end{definition}

\begin{definition}[Star product]
\dcref{ch2:1.1}
\label{topping:notn-star-product}
\source{topping:page-0027}
\group{topping-background}
\lean{Topping.IsStarProduct, Topping.tensorProd, Topping.permSlots,
  Topping.IsStarProduct.sub}
\leanok
The expression $A*B$ denotes any tensor obtained as a universal linear
combination of tensors formed from $A\otimes B$ by metric contractions, index
raising or lowering, and permutations of tensor factors.
\end{definition}

\begin{definition}[Divergence]
\dcref{ch2:1.1}
\label{topping:notn-divergence}
\source{topping:page-0028}
\group{topping-background}
\lean{Topping.divergence, Topping.divergence_apply}
\leanok
For $T\in\Gamma(\otimes^kT^*\M)$, define
\[
\delta(T):=-\tr_{12}\nabla T.
\]
\end{definition}

\begin{definition}[Gravitation tensor]
\dcref{ch2:1.1}
\label{topping:notn-gravitation-tensor}
\source{topping:page-0028}
\group{topping-background}
\lean{Topping.gravitationTensor, Topping.gravitationTensor_apply,
  Topping.trace₂}
\leanok
For a symmetric $2$-tensor $T$, define
\[
G(T):=T-\frac12(\tr T)g.
\]
\end{definition}

\begin{remark}[Background identities]
\dcref{ch2:1.1}
\label{topping:notn-background-remark}
\source{topping:page-0026,topping:page-0027,topping:page-0028}
\group{topping-background}
\lean{Topping.dir_scalarCurvatureAt_eq_two_divRicci,
  Topping.covDerivAlong_metricTensorField_eq_zero,
  Topping.divergence_smul_metricTensorField,
  Topping.divergence_gravitationTensor_ricciTensorField}
The conventions above imply the shorthand identities
\[
R(.,.)A=A*\Rm,\qquad
\nabla(A*B)=(\nabla A)*B+A*(\nabla B),
\]
and the contracted Bianchi identity is
\[
\delta G(\Ric)=\delta\Ric+\frac12dR=0.
\]
\end{remark}

\section{Einstein metrics}

\begin{proposition}[Constancy of the Einstein factor]
\dcref{ch2:2}
\label{topping:einstein-metrics-const-scalar}
\source{topping:page-0029}
\group{topping-background}
\lean{Topping.IsEinsteinWithFactor,
  Topping.scalarCurvatureAt_eq_finrank_mul_of_isEinsteinWithFactor,
  Topping.finrank_sub_two_mul_dirTangent_lam_eq_zero,
  Topping.dirTangent_lam_eq_zero_of_isEinsteinWithFactor,
  Topping.dir_scalarCurvatureAt_eq_zero_of_isEinsteinWithFactor,
  Topping.lam_const_of_isEinsteinWithFactor,
  Topping.scalarCurvatureAt_const_of_isEinsteinWithFactor}
\leanok
If $\Ric(g)=\lambda g$ for a function $\lambda:\M\to\mathbb{R}$ on a connected
$\M$ and $n\ne2$, then the scalar curvature and $\lambda$ are constant. Thus the
constant-factor and function-factor formulations of the Einstein condition
coincide in these dimensions.
\end{proposition}

\section{Deformation of geometric quantities}

Let $g_t$ be a smooth family of metrics and write
$h:=\partial g_t/\partial t$.

\begin{proposition}[Variation of the Levi-Civita connection]
\dcref{ch2:3.1}
\label{topping:variation-levi-civita}
\source{topping:page-0029}
\group{topping-background}
\lean{Topping.IsConnectionVariation, Topping.isConnectionVariation_unique,
  Topping.IsMetricVariationOn, Topping.covTensorOfBilin}
For
\[
\Pi(X,Y):=\frac{\partial}{\partial t}(\nabla_XY),
\]
one has
\[
\langle\Pi(X,Y),Z\rangle
=\frac12\bigl[(\nabla_Yh)(X,Z)+(\nabla_Xh)(Y,Z)-(\nabla_Zh)(X,Y)\bigr].
\]
\end{proposition}

\begin{remark}[Variation of a covariant derivative of a vector field]
\dcref{ch2:3.2}
\label{topping:variation-vector-field}
\source{topping:page-0030}
\group{topping-background}
\lean{Topping.HasVectorFieldVariationOn}
For a time-dependent vector field $V=V(t)$,
\[
\frac{\partial}{\partial t}\nabla_XV
=\Pi(X,V)+\nabla_X\frac{\partial V}{\partial t}.
\]
\end{remark}

\begin{remark}[Variation of covariant derivatives]
\dcref{ch2:3.3}
\label{topping:variation-covariant-derivative-star}
\source{topping:page-0030}
\group{topping-background}
\lean{Topping.HasCovDerivVariationStarOn}
For a fixed tensor field $A$,
\[
\frac{\partial}{\partial t}\nabla A=A*\nabla h,
\]
whereas for a time-dependent tensor field,
\[
\frac{\partial}{\partial t}\nabla A-\nabla\frac{\partial A}{\partial t}=A*\nabla h.
\]
\end{remark}

\begin{proposition}[Variation of the curvature operator]
\dcref{ch2:3.4}
\label{topping:variation-curvature-operator}
\source{topping:page-0030}
\group{topping-background}
\lean{Topping.HasCurvatureOperatorVariationOn,
  Topping.covDerivConnectionVariation}
The curvature operator satisfies
\[
\frac{\partial}{\partial t}R(X,Y)W=(\nabla_Y\Pi)(X,W)-(\nabla_X\Pi)(Y,W).
\]
\end{proposition}

\begin{proposition}[Variation of the Riemann tensor]
\dcref{ch2:3.5}
\label{topping:variation-riemann-tensor}
\source{topping:page-0031}
\group{topping-background}
\lean{Topping.HasRiemannVariationOn}
\[
\begin{aligned}
\frac{\partial}{\partial t}\Rm(X,Y,W,Z)
={}&\frac12\bigl[h(R(X,Y)W,Z)-h(R(X,Y)Z,W)\bigr]\\
&+\frac12\bigl[\nabla^2_{Y,W}h(X,Z)-\nabla^2_{X,W}h(Y,Z)\\
&\qquad+\nabla^2_{X,Z}h(Y,W)-\nabla^2_{Y,Z}h(X,W)\bigr].
\end{aligned}
\]
\end{proposition}

\begin{proposition}[Variation of a trace]
\dcref{ch2:3.6}
\label{topping:variation-trace}
\source{topping:page-0031}
\group{topping-background}
\lean{Topping.HasTraceVariationOn, Topping.bilinTraceAt, Topping.bilinInnerAt,
  Topping.bilinInnerAt_metric}
For $\alpha\in\Gamma(\otimes^2T^*\M)$ depending on $t$,
\[
\frac{\partial}{\partial t}(\tr\alpha)
=-\langle h,\alpha\rangle+\tr\frac{\partial\alpha}{\partial t}.
\]
\end{proposition}

\begin{definition}[Lichnerowicz Laplacian]
\dcref{ch2:3.7}
\label{topping:variation-lichnerowicz}
\source{topping:page-0031}
\group{topping-background}
\lean{Topping.lichnerowiczLaplacian,
  Topping.lichnerowiczLaplacian_ricciTensorField_apply}
\leanok
For a symmetric $2$-tensor $h$,
\[
(\Delta_{\mathcal L}h)(X,W):=(\Delta h)(X,W)-h(X,\Ric(W))-h(W,\Ric(X))
+2\tr h(R(X,\cdot)W,\cdot).
\]
\end{definition}

\begin{proposition}[Variation of the Ricci tensor]
\dcref{ch2:3.7}
\label{topping:variation-ricci}
\source{topping:page-0031,topping:page-0032}
\group{topping-background}
\lean{Topping.HasRicciVariationOn, Topping.HasRicciVariationPointwiseOn}
\[
\frac{\partial}{\partial t}\Ric
=-\frac12\Delta_{\mathcal L}h-\frac12\mathcal L_{(\delta G(h))^\#}g.
\]
Equivalently,
\[
\frac{\partial}{\partial t}\Ric(X,W)
=\Delta\Ric(X,W)-2\langle\Ric(X),\Ric(W)\rangle
+2\langle\Rm(X,\cdot,W,\cdot),\Ric\rangle.
\]
\end{proposition}

\begin{remark}[Lie derivative and Hessian]
\dcref{ch2:3.8}
\label{topping:lie-derivative-symmetric-gradient}
\source{topping:page-0032}
\group{topping-background}
\lean{Topping.symmetricGradient, Topping.symmetricGradient_symm,
  Topping.symmetricGradient_gradientField_eq_two_hessian}
For a $1$-form $\omega$,
\[
\mathcal L_{\omega^\#}g(X,W)=\nabla\omega(X,W)+\nabla\omega(W,X).
\]
In particular, for $\omega=df$,
\[
\mathcal L_{(df)^\#}g=2\Hess(f).
\]
\end{remark}

\begin{proposition}[Variation of scalar curvature]
\dcref{ch2:3.9}
\label{topping:variation-scalar-curvature}
\source{topping:page-0032,topping:page-0033,topping:page-0037}
\group{topping-background}
\lean{Topping.HasScalarVariationOn, Topping.ricciBilinAt,
  Topping.bilinInnerAt_neg_two_ricci, Topping.divergence_ricciTensorField,
  Topping.hasScalarVariationOn_of_isRicciFlowOn}
\[
\frac{\partial R}{\partial t}=-\langle h,\Ric\rangle+\delta^2h-\Delta(\tr h).
\]
\end{proposition}

\begin{proposition}[Variation of the divergence of a one-form]
\dcref{ch2:3.10}
\label{topping:variation-divergence-one-form}
\source{topping:page-0033,topping:page-0038}
\group{topping-background}
\lean{Topping.HasDivergenceOneFormVariationOn,
  Topping.oneFormInnerAt, Topping.covDerivOneFormBilin}
For a time-dependent $1$-form $\omega$,
\[
\frac{\partial}{\partial t}\delta\omega
=\delta\frac{\partial\omega}{\partial t}+\langle h,\nabla\omega\rangle
-\langle\delta G(h),\omega\rangle.
\]
\end{proposition}

\begin{proposition}[Variation of the divergence of the gravitation tensor]
\dcref{ch2:3.11}
\label{topping:variation-divergence-gravitation}
\source{topping:page-0033,topping:page-0039,topping:page-0040}
\group{topping-background}
\lean{Topping.HasDivergenceGravitationVariationOn,
  Topping.oneFormSharp, Topping.oneFormCovec}
If $T\in\Gamma(\Sym^2T^*\M)$ is independent of $t$, then, for every vector
field $Z$,
\[
\left(\frac{\partial}{\partial t}\delta G(T)\right)Z
=-T\bigl((\delta G(h))^\#,Z\bigr)
+\langle h,\nabla T(\cdot,\cdot;Z)-\frac12\nabla_ZT\rangle.
\]
\end{proposition}

\begin{proposition}[Variation of the volume form]
\dcref{ch2:3.12}
\label{topping:variation-volume-form}
\source{topping:page-0033,topping:page-0040}
\group{topping-background}
\lean{Topping.HasVolumeFormVariationOn,
  Topping.selfChartVolumeDensityAt,
  Topping.chartBilinTraceAt_eq_bilinTraceAt,
  Topping.hasVolumeFormVariationOn_selfChart,
  Topping.chartBilinTraceAt_neg_two_ricci_self,
  Topping.hasVolumeFormVariationOn_selfChart_of_isRicciFlowOn,
  Topping.hasDerivWithinAt_selfChartVolumeDensityAt_of_isRicciFlowOn}
\leanok
\[
\frac{\partial}{\partial t}dV=\frac12(\tr h)\,dV.
\]
\end{proposition}

\section{Laplacian of the curvature tensor}

\begin{proposition}[Rough Laplacian of curvature]
\dcref{ch2:4.1}
\label{topping:laplacian-curvature-tensor}
\source{topping:page-0041,topping:page-0042}
\group{topping-background}
\lean{Topping.HasCurvatureLaplacianFormula, Topping.curvatureB,
  Topping.curvatureB_swap_pairs, Topping.curvatureB_swap_within}
\begin{align*}
(\Delta\Rm)(X,Y,W,Z)={}&-\nabla^2_{Y,W}\Ric(X,Z)+\nabla^2_{X,W}\Ric(Y,Z)\\
&-\nabla^2_{X,Z}\Ric(Y,W)+\nabla^2_{Y,Z}\Ric(X,W)\\
&-\Ric(R(W,Z)Y,X)+\Ric(R(W,Z)X,Y)\\
&-2\bigl(B(X,Y,W,Z)-B(X,Y,Z,W)\\
&\qquad+B(X,W,Y,Z)-B(X,Z,Y,W)\bigr).
\end{align*}
\end{proposition}

\section{Evolution under Ricci flow}

\begin{proposition}[Curvature evolution]
\dcref{ch2:5.1}
\label{topping:ricci-flow-curvature-evolution}
\source{topping:page-0042,topping:page-0043}
\group{topping-background}
\lean{Topping.HasCurvatureEvolutionOn, Topping.riemannTensorField,
  Topping.curvatureEvolutionCorrection,
  Topping.isStarProduct_curvatureEvolutionCorrection,
  Topping.HasCurvatureEvolutionComponentsOn,
  Topping.hasCurvatureEvolutionOn_of_components,
  Topping.hasCurvatureEvolutionComponentsOn_of_variation,
  Topping.hasCurvatureEvolutionOn_of_variation,
  Topping.hasCurvatureEvolutionComponentsOn_interior_of_isRicciFlowOn,
  Topping.hasCurvatureEvolutionOn_interior_of_isRicciFlowOn}
Under $\partial_tg=-2\Ric$,
\[
\frac{\partial}{\partial t}\Rm=\Delta\Rm+\Rm*\Rm.
\]
In components,
\begin{align*}
\frac{\partial}{\partial t}\Rm(X,Y,W,Z)={}&(\Delta\Rm)(X,Y,W,Z)\\
&-\Ric(R(X,Y)W,Z)+\Ric(R(X,Y)Z,W)\\
&-\Ric(R(W,Z)X,Y)+\Ric(R(W,Z)Y,X)\\
&+2\bigl(B(X,Y,W,Z)-B(X,Y,Z,W)\\
&\qquad+B(X,W,Y,Z)-B(X,Z,Y,W)\bigr).
\end{align*}
\end{proposition}

\begin{remark}[Heat-type form of curvature evolution]
\dcref{ch2:5.2}
\label{topping:ricci-flow-curvature-heat}
\source{topping:page-0043}
\group{topping-background}
\lean{Topping.hasCurvatureEvolutionOn_iff_heat_type}
The curvature evolution has heat-type form
\[
\frac{\partial}{\partial t}\Rm=\Delta\Rm+\Rm*\Rm.
\]
\end{remark}

\begin{proposition}[Ricci-tensor evolution]
\dcref{ch2:5.3}
\label{topping:ricci-flow-ricci-evolution}
\source{topping:page-0043}
\group{topping-background}
\lean{Topping.HasRicciEvolutionOn, Topping.ricciTensorField,
  Topping.HasRicciVariationPointwiseOn,
  Topping.hasRicciEvolutionOn_iff_pointwise,
  Topping.lichnerowiczLaplacian_ricciTensorField_apply}
Under Ricci flow,
\[
\frac{\partial}{\partial t}\Ric=\Delta_{\mathcal L}(\Ric),
\]
or, pointwise,
\[
\frac{\partial}{\partial t}\Ric(X,W)
=\Delta\Ric(X,W)-2\langle\Ric(X),\Ric(W)\rangle
+2\langle\Rm(X,\cdot,W,\cdot),\Ric\rangle.
\]
\end{proposition}

\begin{proposition}[Scalar-curvature evolution]
\dcref{ch2:5.4}
\label{topping:ricci-flow-scalar-evolution}
\source{topping:page-0043}
\group{topping-background}
\lean{Topping.HasScalarCurvatureEvolutionOn,
  Topping.hasScalarCurvatureEvolutionOn_iff,
  Topping.hasScalarCurvatureEvolutionOn_of_hasScalarVariationOn,
  Topping.hasScalarCurvatureEvolutionOn_of_hasScalarVariationOn',
  Topping.hasScalarCurvatureEvolutionOn_of_isRicciFlowOn}
Under Ricci flow,
\[
\frac{\partial R}{\partial t}=\Delta R+2|\Ric|^2.
\]
\end{proposition}

\begin{remark}[Trace decomposition estimate]
\dcref{ch2:5.4}
\label{topping:ricci-flow-scalar-ineq}
\source{topping:page-0043}
\group{topping-background}
\lean{Topping.traceFreeRicciAt, Topping.ricciTensorAt_eq_traceFreeRicciAt_add,
  Topping.ricciNormSqAt, Topping.scalarCurvatureAt_sq_div_finrank_le_ricciNormSqAt}
\leanok
Writing $\Ric=\mathring{\Ric}+\frac{R}{n}g$ gives
\[
|\Ric|^2\ge\frac{R^2}{n}.
\]
\end{remark}
